% Options for packages loaded elsewhere
% Options for packages loaded elsewhere
\PassOptionsToPackage{unicode}{hyperref}
\PassOptionsToPackage{hyphens}{url}
%
\documentclass[
  english,
  russian,
  12pt,
  a4paper,
  DIV=11,
  numbers=noendperiod]{scrreprt}
\usepackage{xcolor}
\usepackage{amsmath,amssymb}
\setcounter{secnumdepth}{5}
\usepackage{iftex}
\ifPDFTeX
  \usepackage[T1]{fontenc}
  \usepackage[utf8]{inputenc}
  \usepackage{textcomp} % provide euro and other symbols
\else % if luatex or xetex
  \usepackage{unicode-math} % this also loads fontspec
  \defaultfontfeatures{Scale=MatchLowercase}
  \defaultfontfeatures[\rmfamily]{Ligatures=TeX,Scale=1}
\fi
\usepackage{lmodern}
\ifPDFTeX\else
  % xetex/luatex font selection
\fi
% Use upquote if available, for straight quotes in verbatim environments
\IfFileExists{upquote.sty}{\usepackage{upquote}}{}
\IfFileExists{microtype.sty}{% use microtype if available
  \usepackage[]{microtype}
  \UseMicrotypeSet[protrusion]{basicmath} % disable protrusion for tt fonts
}{}
\usepackage{setspace}
% Make \paragraph and \subparagraph free-standing
\makeatletter
\ifx\paragraph\undefined\else
  \let\oldparagraph\paragraph
  \renewcommand{\paragraph}{
    \@ifstar
      \xxxParagraphStar
      \xxxParagraphNoStar
  }
  \newcommand{\xxxParagraphStar}[1]{\oldparagraph*{#1}\mbox{}}
  \newcommand{\xxxParagraphNoStar}[1]{\oldparagraph{#1}\mbox{}}
\fi
\ifx\subparagraph\undefined\else
  \let\oldsubparagraph\subparagraph
  \renewcommand{\subparagraph}{
    \@ifstar
      \xxxSubParagraphStar
      \xxxSubParagraphNoStar
  }
  \newcommand{\xxxSubParagraphStar}[1]{\oldsubparagraph*{#1}\mbox{}}
  \newcommand{\xxxSubParagraphNoStar}[1]{\oldsubparagraph{#1}\mbox{}}
\fi
\makeatother


\usepackage{longtable,booktabs,array}
\usepackage{calc} % for calculating minipage widths
% Correct order of tables after \paragraph or \subparagraph
\usepackage{etoolbox}
\makeatletter
\patchcmd\longtable{\par}{\if@noskipsec\mbox{}\fi\par}{}{}
\makeatother
% Allow footnotes in longtable head/foot
\IfFileExists{footnotehyper.sty}{\usepackage{footnotehyper}}{\usepackage{footnote}}
\makesavenoteenv{longtable}
\usepackage{graphicx}
\makeatletter
\newsavebox\pandoc@box
\newcommand*\pandocbounded[1]{% scales image to fit in text height/width
  \sbox\pandoc@box{#1}%
  \Gscale@div\@tempa{\textheight}{\dimexpr\ht\pandoc@box+\dp\pandoc@box\relax}%
  \Gscale@div\@tempb{\linewidth}{\wd\pandoc@box}%
  \ifdim\@tempb\p@<\@tempa\p@\let\@tempa\@tempb\fi% select the smaller of both
  \ifdim\@tempa\p@<\p@\scalebox{\@tempa}{\usebox\pandoc@box}%
  \else\usebox{\pandoc@box}%
  \fi%
}
% Set default figure placement to htbp
\def\fps@figure{htbp}
\makeatother



\ifLuaTeX
\usepackage[bidi=basic,provide=*]{babel}
\else
\usepackage[bidi=default,provide=*]{babel}
\fi
% get rid of language-specific shorthands (see #6817):
\let\LanguageShortHands\languageshorthands
\def\languageshorthands#1{}


\setlength{\emergencystretch}{3em} % prevent overfull lines

\providecommand{\tightlist}{%
  \setlength{\itemsep}{0pt}\setlength{\parskip}{0pt}}



 
\usepackage[style=gost-numeric,backend=biber,langhook=extras,autolang=other*]{biblatex}
\addbibresource{bib/cite.bib}

\usepackage[]{csquotes}

\usepackage{indentfirst}
\usepackage{float}
\floatplacement{figure}{H}
\IfFileExists{plex-otf.sty}{
  %% Full TeXlive
  \usepackage[
    % math,
    RM={Scale=0.94},SS={Scale=0.94},SScon={Scale=0.94},TT={Scale=MatchLowercase,FakeStretch=0.9},
    DefaultFeatures={Ligatures=Common}
  ]{plex-otf}
  % \usepackage[Scale=MatchUppercase]{juliamono}
}{
  %% TinyTeX
  \usepackage{libertine}
}
\KOMAoption{captions}{tableheading}
\makeatletter
\@ifpackageloaded{caption}{}{\usepackage{caption}}
\AtBeginDocument{%
\ifdefined\contentsname
  \renewcommand*\contentsname{Содержание}
\else
  \newcommand\contentsname{Содержание}
\fi
\ifdefined\listfigurename
  \renewcommand*\listfigurename{Список иллюстраций}
\else
  \newcommand\listfigurename{Список иллюстраций}
\fi
\ifdefined\listtablename
  \renewcommand*\listtablename{Список таблиц}
\else
  \newcommand\listtablename{Список таблиц}
\fi
\ifdefined\figurename
  \renewcommand*\figurename{Рисунок}
\else
  \newcommand\figurename{Рисунок}
\fi
\ifdefined\tablename
  \renewcommand*\tablename{Таблица}
\else
  \newcommand\tablename{Таблица}
\fi
}
\@ifpackageloaded{float}{}{\usepackage{float}}
\floatstyle{ruled}
\@ifundefined{c@chapter}{\newfloat{codelisting}{h}{lop}}{\newfloat{codelisting}{h}{lop}[chapter]}
\floatname{codelisting}{Список}
\newcommand*\listoflistings{\listof{codelisting}{Листинги}}
\makeatother
\makeatletter
\makeatother
\makeatletter
\@ifpackageloaded{caption}{}{\usepackage{caption}}
\@ifpackageloaded{subcaption}{}{\usepackage{subcaption}}
\makeatother
\usepackage{bookmark}
\IfFileExists{xurl.sty}{\usepackage{xurl}}{} % add URL line breaks if available
\urlstyle{same}
\hypersetup{
  pdftitle={Шаблон отчёта по лабораторной работе},
  pdfauthor={Дмитрий Сергеевич Кулябов},
  pdflang={ru-RU},
  hidelinks,
  pdfcreator={LaTeX via pandoc}}


\title{Шаблон отчёта по лабораторной работе}
\usepackage{etoolbox}
\makeatletter
\providecommand{\subtitle}[1]{% add subtitle to \maketitle
  \apptocmd{\@title}{\par {\large #1 \par}}{}{}
}
\makeatother
\subtitle{Простейший вариант}
\author{Дмитрий Сергеевич Кулябов}
\date{}
\begin{document}
\maketitle

\renewcommand*\contentsname{Содержание}
{
\setcounter{tocdepth}{1}
\tableofcontents
}
\listoffigures
\listoftables

\setstretch{1.5}
\chapter{Цель
работы}\label{ux446ux435ux43bux44c-ux440ux430ux431ux43eux442ux44b}

\autocite{tuis}

\chapter{Задание}\label{ux437ux430ux434ux430ux43dux438ux435}

\chapter{Выполнение лабораторной
работы}\label{ux432ux44bux43fux43eux43bux43dux435ux43dux438ux435-ux43bux430ux431ux43eux440ux430ux442ux43eux440ux43dux43eux439-ux440ux430ux431ux43eux442ux44b}

Посмотрим информацию о mc (рис. \autocite*{fig:001}).

\begin{figure}

{\centering \pandocbounded{\includegraphics[keepaspectratio]{image/1.png}}

}

\caption{Информация и mc}

\end{figure}%

И запустим его. Так он выглядит (рис. \autocite*{fig:002}).

\begin{figure}

{\centering \pandocbounded{\includegraphics[keepaspectratio]{image/2.png}}

}

\caption{Интерфейс}

\end{figure}%

С помощью зажатого шифта выделим файлы (рис. \autocite*{fig:003}).

\begin{figure}

{\centering \pandocbounded{\includegraphics[keepaspectratio]{image/3.png}}

}

\caption{Выделение}

\end{figure}%

Через меню Файл скопируем каталог (рис. \autocite*{fig:004}).

\begin{figure}

{\centering \pandocbounded{\includegraphics[keepaspectratio]{image/4.png}}

}

\caption{Копирование}

\end{figure}%

Через меню Файл переместим каталог (рис. \autocite*{fig:005}).

\begin{figure}

{\centering \pandocbounded{\includegraphics[keepaspectratio]{image/5.png}}

}

\caption{Перемещение}

\end{figure}%

Откроем через правую панель информацию (рис. \autocite*{fig:006}).

\begin{figure}

{\centering \pandocbounded{\includegraphics[keepaspectratio]{image/6.png}}

}

\caption{Информация}

\end{figure}%

Откроем через правую панель быстрый просмотр (рис. \autocite*{fig:007}).

\begin{figure}

{\centering \pandocbounded{\includegraphics[keepaspectratio]{image/7.png}}

}

\caption{Быстрый просмотр}

\end{figure}%

Откроем через правую панель древо каталогов (рис. \autocite*{fig:008}).

\begin{figure}

{\centering \pandocbounded{\includegraphics[keepaspectratio]{image/8.png}}

}

\caption{Древо каталогов}

\end{figure}%

С помощью F3 просмотрим содержимое файла (рис. \autocite*{fig:009}).

\begin{figure}

{\centering \pandocbounded{\includegraphics[keepaspectratio]{image/9.png}}

}

\caption{Просмотр файла}

\end{figure}%

С помощью F4 отредактируем файл (рис. \autocite*{fig:010}).

\begin{figure}

{\centering \pandocbounded{\includegraphics[keepaspectratio]{image/10.png}}

}

\caption{Редактирование файла}

\end{figure}%

Создадим каталог через меню Команда (рис. \autocite*{fig:011}).

\begin{figure}

{\centering \pandocbounded{\includegraphics[keepaspectratio]{image/11.png}}

}

\caption{Создание каталога}

\end{figure}%

Скопируем каталог через меню Команда (рис. \autocite*{fig:012}).

\begin{figure}

{\centering \pandocbounded{\includegraphics[keepaspectratio]{image/12.png}}

}

\caption{Копирование файла}

\end{figure}%

Найдём файл .txt, содержащий .d через меню Команда (рис.
\autocite*{fig:013}).

\begin{figure}

{\centering \pandocbounded{\includegraphics[keepaspectratio]{image/13.png}}

}

\caption{Поиск}

\end{figure}%

Откроем через Команда историю, и выполним какую-нибудь команду оттуда
(рис. \autocite*{fig:014}).

\begin{figure}

{\centering \pandocbounded{\includegraphics[keepaspectratio]{image/14.png}}

}

\caption{История}

\end{figure}%

Через Команда откроем каталоги быстрого доступа, где выберем домашний
каталог, чтобы быстро перейти в него (рис. \autocite*{fig:015}).

\begin{figure}

{\centering \pandocbounded{\includegraphics[keepaspectratio]{image/15.png}}

}

\caption{Перемещение в домашний каталог}

\end{figure}%

Откроем файл расширений (рис. \autocite*{fig:016}).

\begin{figure}

{\centering \pandocbounded{\includegraphics[keepaspectratio]{image/16.png}}

}

\caption{Файл расширений}

\end{figure}%

Откроем файл меню (рис. \autocite*{fig:017}).

\begin{figure}

{\centering \pandocbounded{\includegraphics[keepaspectratio]{image/17.png}}

}

\caption{Файл меню}

\end{figure}%

Посмотрим настройки (рис. \autocite*{fig:018}).

\begin{figure}

{\centering \pandocbounded{\includegraphics[keepaspectratio]{image/18.png}}

}

\caption{Просмотр настроек}

\end{figure}%

Создадим файл и вставим в него что-нибудь (рис. \autocite*{fig:019}).

\begin{figure}

{\centering \pandocbounded{\includegraphics[keepaspectratio]{image/19.png}}

}

\caption{Вставка текста в файл}

\end{figure}%

Удалим строку через ctrl + u (рис. \autocite*{fig:020}).

\begin{figure}

{\centering \pandocbounded{\includegraphics[keepaspectratio]{image/20.png}}

}

\caption{Удаление строки}

\end{figure}%

Переместим строку через F6 (рис. \autocite*{fig:021}).

\begin{figure}

{\centering \pandocbounded{\includegraphics[keepaspectratio]{image/21.png}}

}

\caption{Перемещение текста}

\end{figure}%

Скопируем строку через F5 (рис. \autocite*{fig:022}).

\begin{figure}

{\centering \pandocbounded{\includegraphics[keepaspectratio]{image/22.png}}

}

\caption{Копирование текста}

\end{figure}%

Отменим действие через ctrl + u (рис. \autocite*{fig:023}).

\begin{figure}

{\centering \pandocbounded{\includegraphics[keepaspectratio]{image/23.png}}

}

\caption{Отмена действия}

\end{figure}%

Напишем что-нибудь в начале и в конце файла (рис. \autocite*{fig:024}).

\begin{figure}

{\centering \pandocbounded{\includegraphics[keepaspectratio]{image/24.png}}

}

\caption{Редактирование начала и конца файла}

\end{figure}%

И откроем какой-нибудь исходный код, где включим подцветку синтаксиса
через ctrl + s (рис. \autocite*{fig:025}).

\begin{figure}

{\centering \pandocbounded{\includegraphics[keepaspectratio]{image/25.png}}

}

\caption{Включение подцветки синтаксиса}

\end{figure}%

\chapter{Выводы}\label{ux432ux44bux432ux43eux434ux44b}

В результате выполнения лабораторной работы появились навыки работы с mc

\chapter{Контрольные
вопросы}\label{ux43aux43eux43dux442ux440ux43eux43bux44cux43dux44bux435-ux432ux43eux43fux440ux43eux441ux44b}

\begin{enumerate}
\def\labelenumi{\arabic{enumi}.}
\item
  Панели также могут быть настроены для отображения информации в одном
  из двух режимов: \enquote{Информация} или \enquote{Дерево}. В режиме
  \enquote{Информация} на панели отображаются данные о текущем файле и
  файловой системе, на которой расположены активные панели. В режиме
  \enquote{Дерево}, на одной из панелей отображается структура дерева
  каталогов.
\item
  Раздел \enquote{Командная строка оболочки (Shell)} содержит список
  команд и сочетаний клавиш, используемых для ввода и редактирования
  команд в командной строке оболочки. Большинство этих команд
  предназначены для быстрого ввода имен файлов и/или каталогов в
  командную строку (чтобы упростить процесс ввода) или для доступа к
  истории команд. Клавиши редактирования строк ввода применяются как при
  редактировании командной строки, так и в других полях ввода, которые
  появляются в различных запросах программы.
\item
  В меню каждой из панелей (левой или правой) можно выбрать один из
  следующих форматов списка:\\
  Стандартный: отображает список файлов и каталогов с указанием их
  размера и времени последнего изменения.\\
  Ускоренный: позволяет задать количество столбцов, на которые
  разбивается панель при выводе списка имен файлов или каталогов без
  дополнительной информации.\\
  Расширенный: помимо названия файла или каталога выводит дополнительные
  сведения о правах доступа, владельце, группе, размере и времени
  последнего изменения.\\
  Определённый пользователем: предоставляет пользователю возможность
  выбрать, какие дополнительные сведения о файле или каталоге следует
  отобразить.
\item
  Просмотр (F3): Позволяет просмотреть содержимое текущего (или
  выделенного) файла без возможности редактирования.\\
  Просмотр вывода команды (M + !): Запускает команду с параметрами
  (аргумент к текущему выбранному файлу) и позволяет просмотреть её
  вывод.\\
  Правка (F4): Открывает текущий (или выделенный) файл для
  редактирования.\\
  Копирование (F5): Копирует один или несколько файлов или каталогов в
  указанное пользователем место.\\
  Права доступа (Ctrl-x c): Позволяет указать (или изменить) права
  доступа к одному или нескольким файлам или каталогам.\\
  Жёсткая ссылка (Ctrl-x l): Создает жёсткую ссылку на текущий (или
  выделенный) файл.\\
  Символическая ссылка (Ctrl-x s): Создает символическую ссылку на
  текущий (или выделенный) файл.\\
  Владелец/группа (Ctrl-x o): Позволяет задать (или изменить) владельца
  и имя группы для одного или нескольких файлов или каталогов.\\
  Права (расширенные): Позволяет изменить расширенные права доступа и
  владение для одного или нескольких файлов или каталогов.\\
  Переименование (F6): Позволяет переименовать (или переместить) один
  или несколько файлов или каталогов.\\
  Создание каталога (F7): Позволяет создать новый каталог.\\
  Удалить (F8): Позволяет удалить один или несколько файлов или
  каталогов.\\
  Выход (F10): Завершает работу Midnight Commander (mc).
\item
  Дерево каталогов: Отображает структуру каталогов в системе.\\
  Поиск файла: Выполняет поиск файлов с заданными параметрами.\\
  Переставить панели: Меняет местами левую и правую панели.\\
  Сравнить каталоги (Ctrl-x d): Сравнивает содержимое двух каталогов.\\
  Размеры каталогов: Отображает размер и время изменения каталога (по
  умолчанию в MC размер каталога не отображается корректно).\\
  История командной строки: Выводит список ранее выполненных в оболочке
  команд.\\
  Каталоги быстрого доступа (Ctrl-): При вызове выполняется быстрая
  смена текущего каталога на один из заданных в списке.\\
  Восстановление файлов: Позволяет восстановить файлы на файловых
  системах ext2 и ext3.\\
  Редактировать файл расширений: Позволяет определить действия при
  запуске файлов с определенным расширением (например, какое программное
  обеспечение использовать для открытия или редактирования файлов с
  расширением doc или docx).\\
  Редактировать файл меню: Позволяет отредактировать контекстное меню
  пользователя, вызываемое по клавише F2.\\
  Редактировать файл расцветки имён: Позволяет настроить цветовую схему
  для имен файлов в зависимости от их типа.
\item
  Конфигурация: Позволяет настроить параметры работы с панелями.\\
  Внешний вид и Настройки панелей: Определяет элементы (строка меню,
  командная строка, подсказки и т. д.), отображаемые при вызове Midnight
  Commander (MC), а также геометрию расположения панелей и цветовую
  схему.\\
  Биты символов: Задает формат обработки информации локальным
  терминалом.\\
  Подтверждение: Позволяет установить или убрать вывод окна с запросом
  подтверждения действий при операциях удаления и перезаписи файлов, а
  также при выходе из программы.\\
  Распознание клавиш: Диалоговое окно используется для тестирования
  функциональных клавиш, клавиш управления курсором и других
  клавиатурных комбинаций.\\
  Виртуальные ФС: Настройки виртуальной файловой системы, такие как
  тайм-аут, пароль и другие параметры.
\item
  F1: Вызов контекстно-зависимой подсказки.\\
  F2: Вызов пользовательского меню с возможностью создания и/или
  дополнения дополнительных функций.\\
  F3: Просмотр содержимого файла, на который указывает подсветка в
  активной панели (без возможности редактирования).\\
  F4: Вызов встроенного в MC редактора для изменения содержания файла,
  на который указывает подсветка в активной панели.\\
  F5: Копирование одного или нескольких файлов, отмеченных в первой
  (активной) панели, в каталог, отображаемый на второй панели.\\
  F6: Перенос одного или нескольких файлов, отмеченных в первой
  (активной) панели, в каталог, отображаемый на второй панели.\\
  F7: Создание подкаталога в каталоге, отображаемом в активной панели.\\
  F8: Удаление одного или нескольких файлов (или каталогов), отмеченных
  в первой (активной) панели файлов.\\
  F9: Вызов меню MC.\\
  F10: Выход из MC.
\item
  Ctrl-y: Удалить строку.\\
  Ctrl-u: Отмена последней операции.\\
  Ins: Вставка/замена.\\
  F7: Поиск (можно использовать регулярные выражения).\\
  Стрелка вверх + F7: Повтор последней операции поиска.\\
  F4: Замена.\\
  F3 (первое нажатие): Начало выделения, (второе нажатие): Окончание
  выделения.\\
  F5: Копировать выделенный фрагмент.\\
  F6: Переместить выделенный фрагмент.\\
  F8: Удалить выделенный фрагмент.\\
  F2: Записать изменения в файл.\\
  F10: Выйти из редактора.
\item
  Вы можете создавать пользовательские команды для быстрого вызова на
  панели Midnight Commander. Для этого просто введите команду в строке
  ввода (в разделе \enquote{Команда}) и нажмите кнопку
  \enquote{Добавить}. Затем введите информативное имя для этой команды.
  В следующий раз вам будет достаточно выбрать это имя из списка, чтобы
  вызвать команду, вместо того чтобы вводить всю команду заново.
\item
  Панель в Midnight Commander (mc) отображает список файлов текущего
  каталога. Абсолютный путь к этому каталогу отображается в заголовке
  панели. Активная панель выделяется подсветкой в заголовке и одной из
  её строк. Управление панелями осуществляется с помощью определённых
  комбинаций клавиш или пунктов меню в Midnight Commander.
\end{enumerate}

\chapter*{Список
литературы}\label{ux441ux43fux438ux441ux43eux43a-ux43bux438ux442ux435ux440ux430ux442ux443ux440ux44b}
\addcontentsline{toc}{chapter}{Список литературы}

\printbibliography[heading=none]





\end{document}
